\documentclass[../../../include/open-logic-section]{subfiles}

\begin{document}

\olfileid{sth}{card-arithmetic}{fix}
\olsection{د $\aleph$ ثابت ټکي}

په \olref[spine][]{chap} کښې مو ويلي وو چې د تعويض اصل توجيه
غواړي، ځکه دا د سټونو مراتبي جوړښت ډېر لوړوي۔ اوس چې مو د اصلي
عددونو يو څه حساب وکړ، د دې جوړښت د لوړوالي يوه بېلګه وړاندې
کولے شو۔

ښکاره ده چې $0 < \aleph_0$، او $1 < \aleph_1$، او
$2 < \aleph_2$\ldots؛ په رښتيا، د اندازې توپير په هر پړاو کښې
لا \emph{لوېږي}۔ نو دا ګومان راولاړېږي چې
$\kappa< \aleph_\kappa$ به د هر ترتيبي عدد $\kappa$ دپاره وي۔

خو د $\ZFC$ په منلو دا ګومان \emph{ناسم} دے۔ په رښتيا،
\emph{د $\aleph$ ثابت ټکي} هم شته؛ يعنې داسې اصلي عددونه
$\kappa$ چې $\kappa=\aleph_\kappa$ وي۔

\begin{prop}\ollabel{alephfixed}
د $\aleph$ يو ثابت ټکی شته۔
\end{prop}

\begin{proof}
د بازګښت په کارولو لاندې تعريف کړئ:
\begin{align*}
	\kappa_0 &= 0\\
	\kappa_{n+1} &= \aleph_{\kappa_n}\\
	\kappa&= \bigcup_{n < \omega}\kappa_n
\end{align*}
اوس $\kappa$ د
\olref[cardinals][classing]{unioncardinalscardinal} له مخې
اصلي عدد دے۔ خو بيا:
\[
	\kappa= \bigcup_{n < \omega} \kappa_{n+1} = 
	\bigcup_{n < \omega}\aleph_{\kappa_n} = 
	\bigcup_{\alpha < \kappa}\aleph_\alpha = \aleph_\kappa
\]
%	By construction, $\kappa$ is the least cardinal greater than each
%	$\kappa_n$. So $\aleph_\kappa$ is the least cardinal greater than
%	each $\aleph_{\kappa_n}$, and hence greater than each
%	$\kappa_{n+1}$. But equally $\kappa$ is the least cardinal greater
%	than each $\kappa_{n+1} = \aleph_{\kappa_n}$. So $\kappa=
%	\aleph_\kappa$.
\end{proof}

بولوس يو وخت د همدغې، زموږ له خوا جوړې شوې، د $\aleph$
د ثابت ټکي په اړه يوه ليکنه وکړه۔ د ليکنې په سر کښې يې د
$\kappa$ له شتون څخه وروسته وويل:
\begin{quote}
	[د $\kappa$ اندازه] د هغو کسانو له نظره چې مخکې يې د سټونو
	له نظريې سره آشنايي نه وي لرلې، \emph{ډېره لويه} شمېره ده؛
	زما په اند دومره لويه چې د هرې هغې نظريې د رښتياوالي په اړه
	شک پيدا کوي چې يوه ادعا يې دا وي چې لږ تر لږه
	$\kappa$ څيزونه شته۔
	\citep[p.~257]{Boolos2000}
\end{quote}
او د ليکنې په پاے کښې يې دا پوښتنه وکړه:
\begin{quote}
	[ايا] موږ شک کوو چې، که څه هم د کيسې په پيل کښې به داسې
	نه وه، خو تر دې ځايه په رسېدو څرخونه بې‌ځايه تاوېږي او
	موږ نور د کوم واقع شوي شي وصف نه اورو؟
	\citep[p.~268]{Boolos2000}
\end{quote}
که موږ په رښتيا له «هر هغه څه چې واقع دي» هاخوا تللي يو،
نو نېغه پړه به پر تعويض واچوو۔ ځکه همدا اصل زموږ ثبوت ته
لار برابروي۔ که داسې وي، نو ګومان کېږي چې بولوس به اړ شي
خپله هغه څو لسيزې پخوانۍ ادعا بيا وګوري چې تعويض «هېڅ
ناغوښتې» پايله نه لري (وګورئ \olref[replacement][extrinsic]{sec})۔

خو ايا د $\kappa$ شتون دومره بد دے؟ دلته ښايي د رسل
\emph{د ټرېسټرم شېنډي تناقض} په پام کښې نيول مرسته وکړي۔
ټرېسټرم شېنډي خپل ژوند په ورځپاڼه کښې ثبتوي، خو د يوې
ورځې د ثبتولو دپاره يو کال وخت اخلي۔ د هر کال په تېرېدو
لا شاته پاتې کېږي: له يو کال وروسته يې يوازې يوه ورځ
ثبت کړې او ۳۶۴ تېرې ورځې يې نه دي ثبت کړې؛ له دوو کلونو
وروسته يې يوازې دوې ورځې ثبت کړې او ۷۲۸ ورځې نه دي ثبت
کړې؛ له درېيو کلونو وروسته يې يوازې درې ورځې ثبت کړې
او ۱۰۹۲ ورځې نه دي ثبت کړې \dots\footnote{کبيسې کلونه
له پامه غورځوو۔} بيا هم، که ټرېسټرم \emph{تلپاتې} وي،
هره ورځ به ثبت کړے شي، ځکه د خپل ژوند $n$مه ورځ
به په $n$م کال کښې ثبت کړي۔ په دې ډول به «د وخت په پاے
کښې» د هغه ورځپاڼه بشپړه وي۔

اوس ولې دا له دې فکره دومره جلا ده چې $\alpha$ له
$\aleph_\alpha$ څخه کوچنے، بلکې په زياتېدونکي او ډېر
ښکاره ډول کوچنے، وي، تر هغې چې $\kappa$ ته ورسېږو،
او هلته دواړه سره برابر شي او $\kappa
= \aleph_\kappa$ وګرځي؟

دا خبره يوې خوا ته پرېږدو، او که $\ZFC$ ومنو، نو د
ثابت ټکي د جوړښتونو په اړه له يوې بلې په زړه پورې
خبرې سره پای ته رسېږو۔ راتلونکې درې پايلې په وجداني
ډول دا ثابتوي چې يو (ناابتدايي) پړاو شته چې د مراتبي
جوړښت پلنوالے يې له لوړوالي سره برابر دے:

\begin{prop}\ollabel{bethfixed}
د $\beth$ يو ثابت ټکی شته، يعنې داسې $\kappa$ چې
$\kappa=\beth_\kappa$ وي۔
\end{prop}

\begin{proof}
لکه په \olref{alephfixed} کښې، خو «$\beth$» د
«$\aleph$» پر ځاے وکاروئ۔
\end{proof}

\begin{prop}\ollabel{stagesize}
$\card{V_{\omega+\alpha}} = \beth_{\alpha}$۔ که
$\omega \ordtimes \omega \leq \alpha$ وي، نو
$\card{V_\alpha} = \beth_\alpha$۔
\end{prop}

\begin{proof}
لومړۍ ادعا د نامتناهيت هاخوا په ساده استقرا ثابتېږي۔
دويمه ادعا له دې راځي چې که $\omega \ordtimes \omega \leq
\alpha$ وي، نو $\omega + \alpha = \alpha$۔ د دې د
ثابتولو دپاره د ترتيبي حساب هغه حقيقتونه کاروو چې په
\olref[ord-arithmetic][]{chap} کښې راغلي دي۔ لومړے پام
وکړئ چې $\omega \ordtimes \omega = \omega \ordtimes (1
\ordplus \omega) = ( \omega  \ordtimes 1) \ordplus
(\omega\ordtimes\omega) = \omega
\ordplus (\omega \ordtimes \omega)$۔ اوس که
$\omega \ordtimes \omega \leq \alpha$ وي، يعنې
$\alpha = (\omega\ordtimes\omega) \ordplus \beta$
د کوم $\beta$ دپاره وي، نو $\omega \ordplus \alpha = \omega
\ordplus ((\omega \ordtimes \omega) \ordplus \beta) =
(\omega \ordplus (\omega \ordtimes \omega)) \ordplus
\beta = (\omega \ordtimes \omega) \ordplus \beta =
\alpha$۔
\end{proof}

\begin{cor}
داسې $\kappa$ شته چې $\card{V_\kappa} = \kappa$۔
\end{cor}

\begin{proof}
هغه $\kappa$ واخلئ چې د \olref{bethfixed} له مخې
د $\beth$ ثابت ټکی دے۔ څرګنده ده چې $\omega \ordtimes
\omega < \kappa$۔ نو د \olref{stagesize} له مخې
$\card{V_\kappa} = \beth_\kappa= \kappa$۔
\end{proof}

د $V_\kappa$ لاندې دومره پړاوونه شته، څومره چې په
$V_\kappa$ کښې !!{element}s شته۔ نو په وجداني توګه
$V_\kappa$ هماغومره پلن دے څومره چې لوړ دے۔
دا د ټرېسټرم شېنډي د کيسې په شان ده: له يوه پړاو څخه
بل ته د \emph{توان سټونو} په اخيستلو ځو، او په هر
ګام کښې مو مراتبي جوړښت \emph{ډېر} لوېږي۔ خو «په
پاے کښې»، يعنې د $\kappa$ په پړاو کښې، د جوړښت
پلنوالے له لوړوالي سره برابر شي۔

ښايي څوک وپوښتي: \emph{د مراتبي جوړښت پلنوالے
څو ځله له لوړوالي سره برابرېږي؟} ځواب دا دے:
\emph{د ترتيبي عددونو هومره ځله۔} خو دا لږې
تشريح ته اړتيا لري۔

يو ترم $\tau$ په لاندې ډول تعريفوو۔ د هر $A$ دپاره:
\begin{align*}
	\tau_0(A) & \defis \card{A}\\
	\tau_{n+1}(A) & \defis \beth_{\tau_n(A)}\\
	\tau(A) & \defis \bigcup_{n < \omega}\tau_n(A)
\intertext{لکه په \olref{bethfixed} کښې، $\tau(A)$ د
$\beth$ ثابت ټکی دے، د هر $A$ دپاره؛ چاپي متن دا
هم «څرګنده» بولي چې $\card{A} < \tau(A)$۔
نو اوس دا بازګښتي تعريف وګورئ:}
	W_0 &\defis 0\\
	W_{\alpha + 1} & \defis \tau(W_\alpha)\\
	W_\alpha & \defis \bigcup_{\beta < \alpha} W_\beta
	\text{، چې $\alpha$ حدي وي}
\end{align*}
\textbf{د سرچينې يادښت (OLSTH-022):} د چاپي «تل لوے»
ادعا ناسمه ده: که د لومړني سټ اصلي شمېر لا د بېت ثابت ټکی
وي، دا تکرار هماغه اصلي شمېر بېرته ورکوي۔ په ځانګړي
ډول د صفر څخه جوړ شوے لومړے بېت ثابت ټکی په بل تالي
پړاو کښې نه لوېږي؛ نو لاندې يو پر يو توب هم په چاپي
تعريف نه ثابتېږي۔ فورمولونه هماغسې پاتې دي او تشه
لا نه ده رغول شوې۔
دا جوړښت د ټولو ترتيبي عددونو دپاره تعريف شوے دے۔
نو په وجداني توګه، $W$ د ترتيبي عددونو څخه د
$\beth$ ثابت ټکو ته «!!a{injection}» ده۔ او لکه
مخکې، $V_{W_\alpha}$ د هر $\alpha$ دپاره له خپل
لوړوالي سره هم‌پلن دے۔

\end{document}
